%% Font size %%
\documentclass[11pt]{article}

%% Load the custom package
\usepackage{Mathdoc}

%% Numéro de séquence %% Titre de la séquence %%
\renewcommand{\centerhead}{S14- Probabilités - Exercices Divers}

%% Spacing commands %%
\renewcommand{\baselinestretch}{1} \setlength{\parindent}{0pt}

%%\pagestyle{empty}

\begin{document}

\begin{exercice}[2][Calculer une probabilité]

On lance un dé à 6 faces.

Calculer la probabilité des événements suivants.

\begin{enumerate}
    \item Obtenir un 3.
    \item Obtenir un nombre pair.
    \item Obtenir un nombre inférieur à 5.
    \item Obtenir un nombre strictement supérieur à 2.
\end{enumerate}

\end{exercice}


\begin{exercice}[2][Boules dans un sac]

Un sac contient 5 boules rouges, 3 boules bleues et 2 boules vertes.

On tire une boule au hasard.

\begin{enumerate}
    \item Combien y a-t-il de boules au total ?
    \item Calculer la probabilité d'obtenir :
    \begin{enumerate}
        \item une boule rouge ;
        \item une boule bleue ;
        \item une boule verte.
    \end{enumerate}
\end{enumerate}

\end{exercice}


\begin{exercice}[2][Cartes numérotées]

On place dans une boîte des cartes numérotées de 1 à 10.

On tire une carte au hasard.

Calculer la probabilité des événements suivants.

\begin{enumerate}
    \item Obtenir un nombre pair.
    \item Obtenir un multiple de 3.
    \item Obtenir un nombre supérieur ou égal à 8.
    \item Obtenir un nombre inférieur à 11.
\end{enumerate}

\end{exercice}


\begin{exercice}[2][Roue de loterie]

Une roue est partagée en 8 parts identiques :
3 rouges, 2 bleues, 2 vertes et 1 jaune.

On fait tourner la roue.

\begin{enumerate}
    \item Calculer la probabilité d'obtenir la couleur rouge.
    \item Calculer la probabilité d'obtenir la couleur verte.
    \item Calculer la probabilité de ne pas obtenir la couleur jaune.
\end{enumerate}

\end{exercice}


\begin{exercice}[2][Vrai ou faux]

Dire si les affirmations suivantes sont vraies ou fausses.

\begin{enumerate}
    \item Une probabilité peut être supérieure à 1.
    \item Une probabilité peut être égale à 0.
    \item Une probabilité peut être négative.
    \item La somme des probabilités de toutes les issues vaut 1.
\end{enumerate}

\end{exercice}


\end{document}
